\section{Schlüsselbegriffe}

\keywords{Internet of Things (IoT), Intrusion Detection System (IDS), Machine Learning, Gradient Boosting, CICIoT2023, Balanced Accuracy, Klassenungleichgewicht, Explainable Artificial Intelligence (XAI), SHAP (SHapley Additive Explanations), Principal Component Analysis (PCA), Cyberangriffserkennung, Feature Importance, Netzwerksicherheit, Vertrauenswürdige Systeme, Anomalieerkennung}

\section{Abkürzungsverzeichnis}\label{abkuerzungen}

%%%%%%%%%%%%%%%%%%%%%%%%%%%%%%%%%%%%%%%%%%%%%%%%%%%%%%%%%%%%%%%%%%%%%%%%%%%%%%%%%%%%%%%%%
%                             Mit Einrückung im Verzeichnis                             %                     
%  Alle Abkürzungen müssen mit: \acro{acronym}[shortname]{fullname} eingetragen werden. %
%                         Im Fließtext mit \ac{acronym} aufrufen.                       %
%                 Erster Aufruf: fullname (shortname), danach shortname.                %                                   
%%%%%%%%%%%%%%%%%%%%%%%%%%%%%%%%%%%%%%%%%%%%%%%%%%%%%%%%%%%%%%%%%%%%%%%%%%%%%%%%%%%%%%%%%

\begin{addmargin}[1.5em]{0em}

    \begin{acronym}[CICIoT2023] %längster shortname muss eingetragen werden.

        \acro{csv}[CSV]{Comma-Separated Values}

        \acro{ddos}[DDoS]{Distributed Denial of Service}

        \acro{dos}[DoS]{Denial of Service}

        \acro{gb}[GB]{Gradient Boosting}

        \acro{https}[HTTPS]{Hypertext Transfer Protocol Secure}

        \acro{ids}[IDS]{Intrusion Detection System}

        \acro{iot}[IoT]{Internet of Things}

        \acro{ivs}[IVS]{Vertrauenswürdige Systeme}

        \acro{ml}[ML]{Machine Learning}

        \acro{pca}[PCA]{Principal Component Analysis}

        \acro{shap}[SHAP]{SHapley Additive Explanations}

        \acro{xai}[XAI]{Explainable Artificial Intelligence}

    \end{acronym}

\end{addmargin}